\section{Results}
\label{sec:results}

\subsection{Dataset and Experimental Setup}

The TON\,IoT Network Dataset~\cite{ref_ton_iot} provides 22.3~million network flow
records across 10 attack categories. A representative stratified sample of $536{,}052$
flows was drawn with class-proportional sampling, preserving all $1{,}052$ MitM instances
and the complete $72{,}805$ ransomware subset (Table~\ref{tab:eda}).
Per-class chronological splits (60/20/20) were applied to prevent temporal data leakage,
yielding $321{,}631$ training, $107{,}210$ validation, and $107{,}211$ test flows.
Stealthy attack classes (scanning, mitm, backdoor, ransomware) collectively
represent $25.38\%$ of the sampled dataset, with mitm constituting
$0.20\%$ of total flows.

\subsection{Stage-1 Anomaly Detection}
\label{sec:results_s1}

Three one-class detectors were trained exclusively on normal traffic
($n_\text{train}^\text{normal}=48{,}000$). Table~\ref{tab:stage1} reports performance
on the validation set with threshold $\theta$ calibrated for Recall~$\geq 0.90$.

\textbf{OCSVM attains the best operational trade-off}, achieving Recall$=0.900$ at
FPR$=0.0149$ --- the only model satisfying the constraint FPR~$\leq 0.15$ --- with
AUC-ROC$=0.9767$ and AUC-PR$=0.9930$.
Isolation Forest achieves comparable recall but substantially higher FPR ($0.2049$),
which would generate excessive false alarms in production ICS environments.
LOF attains intermediate performance (FPR$=0.1632$, AUC-ROC$=0.9279$).

OCSVM flags $245{,}090/321{,}631$ training flows ($76.2\%$) and $82{,}296/107{,}211$
test flows ($76.8\%$) as anomalous, reflecting the dataset's elevated attack-to-normal
ratio.

\subsection{Stage-2 Classification and Hybrid Pipeline}
\label{sec:results_s2}

Table~\ref{tab:stage2_main} reports performance for three standalone classifiers
and three hybrid pipelines (OCSVM + classifier + ADASYN).

\textbf{Standalone classifiers.}
XGBoost achieves the highest F1-macro ($0.8673$), MCC ($0.9236$), and AUC-PR ($0.9435$)
among standalone configurations, establishing it as the primary baseline.
XGBoost's backdoor recall is $0.3615$, correctly classifying fewer than half
of backdoor events.
Random Forest ranks second (F1-macro$=0.7108$), while LightGBM underperforms
substantially (F1-macro$=0.4141$) due to an incompatibility between the
\texttt{class\_weight='balanced'} parameter and \texttt{numpy} array inputs in the
evaluated library version~\cite{ref_lgb}; this constitutes a methodological pitfall
documented in Section~\ref{sec:discussion_limits}.

\textbf{Hybrid pipelines.}
The OCSVM + XGBoost configuration achieves a \textbf{backdoor recall of $0.9250$},
compared to $0.3615$ for standalone XGBoost --- a gain of $+0.5635$ ($+56.4\%$).
This security-relevant improvement comes at the cost of F1-macro$=0.7337$
($-13.36\%$ vs.\ standalone XGBoost).
The Friedman test ($\chi^2=45.66$, $p<0.0001$) confirms global significance.
All Wilcoxon pairwise tests (OCSVM+XGB vs.\ each standalone) are significant
after Bonferroni correction ($p_\text{Bonf}=0.018$, $\alpha=0.05$).

\subsection{Ablation Study}
\label{sec:results_ablation}

Table~\ref{tab:ablation} decomposes each pipeline component's contribution.

Configuration~A1 (XGBoost + anomaly score) improves marginally over baseline
($+0.0031$ F1-macro), confirming that the Stage-1 anomaly score carries limited
discriminative information beyond the original features.
Configuration~A2 (XGBoost + ADASYN, no filtering) \textbf{degrades} F1-macro by
$-0.0811$ relative to baseline. This unexpected result indicates that ADASYN's
synthetic oversampling of rare classes (mitm, ransomware) introduces
distributional artifacts that confound XGBoost's decision boundaries.
Configuration~A3 (Stage-1 filter alone) incurs the largest penalty
($-0.1815$ F1-macro), as filtering without oversampling leaves the flagged
subset severely imbalanced.
Configurations~A4 and A5 recover to F1-macro$=0.7373$ ($-0.0772$ vs. A0)
while achieving backdoor recall$=0.925$ --- demonstrating that the hybrid's
contribution is targeted class-specific recall improvement rather than global
performance gain.

\subsection{Hypotheses Evaluation}
\label{sec:results_hypotheses}

\textbf{H1} (\textit{Hybrid improves F1-macro on stealthy classes}) is
\textbf{REJECTED on aggregate F1-macro}: the best hybrid configuration
(OCSVM~+~XGBoost, F1-macro$=0.7337$) does not surpass standalone XGBoost
($0.8673$, $\Delta=-0.1336$).
However, H1 is \textbf{VALIDATED on backdoor recall}: the hybrid achieves
$0.9250$ vs.\ $0.3615$ for standalone XGBoost ($+56.4\%$), a statistically
significant and operationally meaningful improvement for ICS security.

\textbf{H2} (\textit{Enrichment factor $\geq 5\times$}) is
\textbf{PARTIALLY VALIDATED}: OCSVM achieves an enrichment of $\times 1.22$
(below the $\times 5$ target), increasing stealthy-class representation
from $25.38\%$ to $31.03\%$. The limited enrichment reflects the dataset's
already-elevated attack proportion (85\% attacks in training data),
which constrains the filter's discriminating power.
